\section{Multi-Robot Coordination}
\label{sec:multi}

Multi-robot coordination involves task allocation, information exchange, state fusion, and motion interaction. Since teams exchange frontiers, maps, observations, poses, or position histories, communication degradation turns local operation and shared-state consistency into separate objectives. Table~\ref{S-tab:multi-robot} in the supplementary material lists representative mechanisms with their evaluation settings.

\subsection{Task Allocation and Spatial Responsibility}

Allocation reduces redundant work through a common state. Centralized planners partition merged maps into dynamic Voronoi cells in CURE~\cite{P_M17}, assign known targets over a sparse Delaunay graph with priority yielding in DMT-CPP~\cite{P_E35}, or assign tasks before the mission while penalizing capability shortfall~\cite{P_M32}. Distributed methods reduce negotiation without removing shared state. CQLite exchanges selected Q-values and frontiers~\cite{P_M25}, LPFE predicts teammates' frontier choices from a common map and last known positions~\cite{P_E34}, and 4CNet predicts map structure from shared trajectories with an estimated confidence~\cite{P_E33}. Such inference avoids messages only while robots hold compatible inputs. URGE abstracts the map into subregions with exploration states and allocates them by jointly minimizing path cost and balancing load~\cite{P_R22}. Its full-set exchange with an explored-overrides-unexplored rule lets regional maps converge after a link returns. Its baseline comparisons assume ideal communication in simulation, whereas three-robot runs in a forest, a parking lot, and a tunnel use an ad hoc mesh network. Terrain-aware partitioning (TAPM-IVP) extends spatial responsibility to uneven ground with communication-radius-constrained regions and traversability-aware subgoals~\cite{P_F09}. In ten simulated trials of Scene~1, each robot succeeds in 90\% of trials, compared with 40\% and 30\% for an ablation that removes both partitioning and greedy subgoal selection. These per-robot rates do not give the probability that the whole team completes the mission. The plane-fitting covariance, moreover, describes terrain scatter rather than propagated pose uncertainty.

\subsection{Compact Sharing and Distributed Estimation}

Compact representations trade bandwidth against what teammates can infer. MR-TopoMap shares topological vertices with place descriptors while keeping local occupancy for avoidance~\cite{P_M36}, Cross-rank exchanges position histories and assigns directions at intersections~\cite{P_F14}, and WiSER-X infers relative positions from radio pings that require an interpretation when they are missing~\cite{P_E32}. SlideSLAM buffers exchanged object observations until a loop closure fixes the frame transform~\cite{P_E31}, whereas Commerge minimizes transmitted scans for central submap registration without adapting to radio capacity~\cite{P_E36}. Bandwidth-efficient collaborative SLAM separates lossy mutual-information keyframe selection from lossless entropy coding under a central back end~\cite{P_S29}. Feedback-guided hierarchical transmission sends a pose-and-factor skeleton, prioritized geometry, and semantic labels in layers and reallocates bandwidth each round from back-end loop-closure and risk signals~\cite{P_R21}. However, its dataset-replay evaluation reports estimated rather than measured collision risk, while the authors note that round-wise guidance can become stale under asynchronous links. PROMO-Orchard gates submap merging with orchard priors and posterior checks against a reference built by the same backbone~\cite{P_S23}.

Distributed estimators must also handle delayed and repeated information. iMESA keeps local incremental factor graphs and updates consensus variables through a two-stage exchange that must complete once started~\cite{P_E40}. riMESA, a direct extension by the same authors, removes this blocking requirement~\cite{P_R23}. It formulates collaborative SLAM as edge-based consensus ADMM, rejects outliers with a robust local solver, and discards failed exchanges, with experiments that replay recorded LiDAR factors under Wi-Fi or radio connectivity models. Multi-Proxy combines lightweight descriptors for progressive loop detection with asynchronous ADMM pose-graph optimization~\cite{P_R20}. Under simulated packet loss and latency, a two-stage Gauss--Seidel baseline fails to converge while Multi-Proxy continues, although the method relies on stable front-end odometry. Parallel graph inference separates current-state from historical-message inference~\cite{P_E37}, whereas CoVOR-SLAM shares poses, covariances, and ranges instead of features~\cite{P_F15}. Swarm-SLAM runs decentralized inference with sparse loop prioritization online on a heterogeneous physical team~\cite{P_S10}, whereas foundation-model loop closing anchors monocular scale with metric odometry in Jetson replay~\cite{P_S11}. Radio geometry adds constraints with their own assumptions, as in cooperative range-only SLAM analyzed offline with all landmarks visible~\cite{P_S27} and formation-controlled UWB anchors whose geometry competes with obstacle clearance~\cite{P_S28}. A prior-graph extension adds informative loop edges to multi-robot coverage paths~\cite{P_N07}, building on single-robot SLAM-aware planning~\cite{P_N05}, although its approximation bound concerns the graph objective rather than physical loop correspondences.

\subsection{Rendezvous, Disconnection, and Conditional Repair}

Communication can become part of the motion objective. A rendezvous-aware explorer creates backtracking frontiers from decaying trajectory records~\cite{P_E38}, whereas a communication-constrained controller schedules sub-team meetings that require full-team agreement to update~\cite{P_E39}. CoCoPlan treats team-wise communication events as decision variables and co-optimizes them with task assignment through branch-and-bound search~\cite{P_R19}. In simulations with up to 100 robots, it reports a 22.4\% higher task-completion rate and 58.6\% lower communication overhead than the compared protocols. MEF-Explore derives rendezvous from entropy fields while assuming that low-rate position exchange is always available~\cite{P_R24}. With three TurtleBot2 robots, it explores 21.32\% faster with a 16.67\% higher success rate than the baseline. All of these methods depend on reaching an encounter, exchanging the required state, and executing the subsequent plan.

During disconnection, robots can continue locally and reconcile later. ROAM continues planning from local maps during observed Wi-Fi cutoffs at the cost of revisits~\cite{P_X08}. Correlation-aware collaboration exchanges states, covariances, and cross-covariances of common landmarks before Kullback--Leibler averaging and Schmidt--Kalman updates, a procedure that excludes faulty observations and corrects uncertain landmarks~\cite{P_F03}. Failure-aware target tracking fuses target estimates by covariance intersection after reconnection and lets a permanently isolated drone signal a hazard~\cite{P_F12}. The communication-aware CBF controller relaxes conflicting constraints and drops unreliable links, with recovery evaluated under partial connectivity~\cite{P_F13}. Planned synchronization makes observation an action, scheduling check-ins in which a TurtleBot pair stops and pans to observe each other~\cite{P_S25}. A range-based integrity monitor reconstructs shared localization errors with distributed ADMM, although its demonstration runs centrally on emulated ranges~\cite{P_N11}. These mechanisms address target estimates, landmarks, links, or observations, although none of them supplies a complete protocol for accepting stale or biased map content after an uncontrolled partition.

\subsection{Motion Interaction and Persistent Service}

Coordination also occurs below the allocation layer. RL-RVO learns a distributed velocity policy from reciprocal velocity-obstacle features with external localization~\cite{P_M38}, predefined-time formation control compensates modeled actuator faults over a fixed graph~\cite{P_X13}, and compensation-based MPC estimates model discrepancy online for four motion-captured nonholonomic robots, leaving deadlock to future work~\cite{P_R17}. Reachability-based platoon control preserves separation under bounded inputs while assuming reliable communication~\cite{P_S04}, whereas distributed barrier control adds a coordinated deadlock response on physical vehicles~\cite{P_S06}. In mEclares, a scheduler reserves return and landing opportunities from energy and timing conditions, with agents falling back to committed paths after disconnection~\cite{P_E57}. These results establish implemented coordination responses, whereas degraded-network validity requires its own fault model and trials.

\subsection{Synthesis}
\label{sec:multi-synthesis}

Coordination methods reduce redundant work through allocation, compact representations, or distributed inference. Lower message volume changes which information is available without resolving inconsistencies between local estimates. Rendezvous planning, asynchronous optimization, and local planning sustain activity during disconnection, whereas reconnection additionally requires validating correspondences, accounting for correlations, and deciding which updates to accept. Evaluations with moving wheeled teams should therefore relate communication loss and reconnection to estimate divergence, coverage, interventions, and recovery time.
